\exercisesfor{7}

\begin{enumerate}
\def\labelenumi{\arabic{enumi}.}
\tightlist
\item
  令 \(\mathfrak{g}\) 是 \(\mathbf{K}\) 上满足下式的完备赋范李代数：
\end{enumerate}

\[
\| [ x, y ] \| \leqslant \| x \| \| y \|
\]

当 \(x, y\) 属于 \(\mathfrak{g}\) 时成立。令 \(\Theta\) 表示所有 \(x \in \mathfrak{g}\) 中满足 \(\| x \| < \frac{1}{3} \log \frac{3}{2}\) 者所成的集合。对于 \(x, y\) 属于 \(\Theta\)，有 \(h(x, y) \in \Theta\)，参见~第 2 号。

\begin{enumerate}
\def\labelenumi{(\alph{enumi})}
\item
  记 \(f(\mathrm{T}) = (1 - \varepsilon^{-\mathrm{T}}) / \mathrm{T}\)，\(g(\mathrm{T}) = 1 / f(\mathrm{T})\)。级数 \(f\) 在整个复平面上收敛，而级数 \(g\) 在半径为 \(2\pi\) 的开圆盘上收敛（参见~《实变函数》，第六章，§2，第 3 号）。推出若 \(z \in \Theta\)，则 \(f(\mathrm{ad}z)\) 和 \(g(\mathrm{ad}z)\) 有定义，且它们是 \(\mathcal{L}(\mathfrak{g};\mathfrak{g})\) 的元素，并互为逆元。
\item
  令 \(\mathrm{D}_2h(x,y)\) 为 \(h\) 在点 \((x,y)\)（该点属于 \(\Theta \times \Theta\)）处的第二个偏导数（《微分与解析流形》，R，1.6.2）。它是 \(\mathcal{L}(\mathfrak{g};\mathfrak{g})\) 的元素。证明
\end{enumerate}

\[
\mathrm{D} _ {2} h (x, y) = g (\operatorname{ad} h (x, y)) \circ f (\operatorname{ad} y).
\]

（使用§6 的习题 3(e)，证明当 \(x\)、\(y\) 充分接近零时该公式成立，再通过解析延拓将其推广到一般情形。）证明公式

\[
f (\operatorname{ad} h (x, y)) \circ \mathrm{D} _ {2} h (x, y) = f (\operatorname{ad} y)
\]

在形式幂级数 \(\tilde{H}\) 的绝对收敛域的每一点都成立（参见~命题 1）。

